\section{Part 2 --- $\theta$ as an Optimization Metric}
\label{sec:deltatheta}

% The engineering arc: use the frozen instrument to measure interventions.
Part~1 produced a fixed, calibrated instrument. We now use it to ask an
engineering question: \emph{when a system-prompt intervention appears to
make a model ``better'' at compliance, what exactly improved?} Because
the bank is Rasch-calibrated, we can read each configuration's ability
$\theta$ on a common scale and compare interventions as differences
$\Delta\theta$.

\subsection{Setup}
We hold the bank fixed at the 1{,}284 frozen items and score nine subject
configurations: three models
(\texttt{llama-3.3-70b-versatile}, \texttt{openai/gpt-oss-20b},
\texttt{llama-3.1-8b-instant}) at a single temperature ($t = 0.4$) under
three system-prompt strictness levels (\texttt{none}, \texttt{neutral},
\texttt{strict}). Restricting the bank to the healthy subset shifts the
baseline's ability, so we re-anchor inside the analysis so that absolute
$\theta$ values remain honest; $\Delta\theta$ comparisons within a model
are invariant to this shift.

\subsection{The dual-judge decomposition}
The strict judge requires section-level citations, so a measured
``prompt effect'' can mean two very different things: the prompt taught
the model to \emph{cite} (a format unlock), or it taught the model to
\emph{reason} better about compliance. To separate them, we re-grade the
same saved responses with the lenient judge (conclusion + reasoning, no
citation) and define, per model,
\begin{align}
\text{reasoning gain} &= \Delta\theta_{\text{lenient}}
  \quad(\text{none}\!\to\!\text{strict}),\\
\text{format unlock}  &= \Delta\theta_{\text{strict}}
  - \Delta\theta_{\text{lenient}}.
\end{align}
The re-grade is offline---it reuses the response text already logged in
Phase~1---so the decomposition adds judge calls but no new solver calls.

\subsection{Results}
Under the strict judge, system-prompt strictness moves all three models,
most dramatically the 70b (Table~\ref{tab:deltatheta},
Figure~\ref{fig:deltatheta}). But the lenient judge flattens the effect
for everything except the 70b: for \texttt{gpt-oss-20b} and the 8b, the
lenient $\Delta\theta$ is essentially zero
(Figure~\ref{fig:decompose}). The honest reading is that
\textbf{system-prompt strictness primarily teaches citation format, not
compliance reasoning}---only the 70b shows a non-trivial reasoning gain
($\approx +2.5$ logits) on top of the format-unlock component.

\begin{table}[t]
\centering
\small
\begin{tabular}{@{}lcccc@{}}
\toprule
model & strict $\Delta\theta$ & lenient $\Delta\theta$
      & format unlock & reasoning gain \\
      & (none$\to$strict) & & & \\
\midrule
\texttt{llama-3.3-70b} & $+5.41$ & $+2.52$ & $+2.89$ & $+2.52$ \\
\texttt{gpt-oss-20b}   & $+0.66$ & $-0.23$ & $+0.89$ & $\approx 0$ \\
\texttt{llama-3.1-8b}  & $+1.69$ & $+0.18$ & $+1.52$ & $\approx 0$ \\
\bottomrule
\end{tabular}
\caption{Ability shift from system-prompt strictness, decomposed.
Strict $\Delta\theta$ is the headline effect; the lenient re-grade
isolates the reasoning component.}
\label{tab:deltatheta}
\end{table}

\begin{figure}[t]
\centering
\includegraphics[width=0.82\linewidth]{fig3_delta_theta.pdf}
\caption{$\Delta\theta$ relative to the \texttt{none} prompt under the
strict judge, by model and strictness level.}
\label{fig:deltatheta}
\end{figure}

\begin{figure}[t]
\centering
\includegraphics[width=0.86\linewidth]{fig4_decompose.pdf}
\caption{Decomposition of the strict-prompt $\Delta\theta$ into a
reasoning gain (lenient $\Delta\theta$) and a format-unlock component
(strict $-$ lenient). Only \texttt{llama-3.3-70b} shows a substantial
reasoning component.}
\label{fig:decompose}
\end{figure}
